%========================================================
%==  CAPÍTULO 2 - FUNDAMENTAÇÃO TEÓRICA
%========================================================
\chapter[Fundamentação Teórica]{Fundamentação Teórica}\label{ch:fundamentacao}

<Apresente os conceitos, técnicas e tecnologias necessários para a
compreensão do trabalho.>

\section{<Primeiro Conceito>}\label{sec:conceito1}

<Desenvolva o conceito, apoiando-se em referências \cite{exemplo-artigo}.>

% Exemplo de inclusão de figura. Coloque o arquivo em imagens/.
%\begin{figure}[!ht]
%    \centering
%    \includegraphics[width=0.75\textwidth]{imagens/exemplo.png}
%    \caption{<Legenda da figura>.}
%    \label{fig:exemplo}
%    \fonte{<Autoria própria ou referência>.}
%\end{figure}

\section{<Segundo Conceito>}\label{sec:conceito2}

<Desenvolva o conceito.>

% Exemplo de tabela.
%\begin{table}[!ht]
%\centering
%\caption{<Legenda da tabela>.}
%\label{tab:exemplo}
%\begin{tabular}{l|c|r}
%\hline
%\textbf{Coluna A} & \textbf{Coluna B} & \textbf{Coluna C} \\
%\hline
%<valor> & <valor> & <valor> \\
%\hline
%\end{tabular}
%\fonte{<Autoria própria ou referência>.}
%\end{table}

\section{Trabalhos Relacionados}\label{sec:relacionados}

<Apresente os trabalhos (artigos, outros TCCs, ferramentas ou sistemas
existentes) mais próximos ao seu e explicite o que diferencia a sua proposta.>
